% Threat Model

\section{Threat Model}
\label{sec:threat-model}

This section defines the adversary capabilities we consider, the security goals \textsc{WasmSentinel} aims to achieve, and the boundaries of our trusted computing base.

\subsection{Adversary Model}

We consider two classes of adversaries:

\textbf{External Attackers.}
These adversaries attempt to compromise the monitored system through various attack vectors:
\begin{itemize}
    \item Execute arbitrary code in user space with unprivileged access
    \item Perform privilege escalation through kernel vulnerabilities or misconfigurations
    \item Exfiltrate sensitive data through covert channels
    \item Deploy persistence mechanisms (rootkits, backdoors)
    \item Evade detection through syscall obfuscation or mimicry attacks
\end{itemize}

\textbf{Malicious Rule Authors.}
A key differentiator of \textsc{WasmSentinel} is its ability to safely execute third-party detection rules.
We consider adversaries who:
\begin{itemize}
    \item Submit malicious detection rules disguised as legitimate threat signatures
    \item Attempt to escape the Wasm sandbox to access host resources
    \item Try to exhaust system resources (CPU, memory) through denial-of-service
    \item Craft adversarial inputs to evade ML-based detection
\end{itemize}

\textbf{Out of Scope.}
We do not consider:
\begin{itemize}
    \item Attackers with root privileges on the host system
    \item Physical attacks or hardware-level compromises
    \item Side-channel attacks against the Wasm runtime (addressed by orthogonal techniques)
    \item Attacks on the eBPF verifier itself (trusted kernel component)
\end{itemize}

\subsection{Security Goals}

\textsc{WasmSentinel} aims to achieve the following security properties:

\begin{table}[t]
\centering
\caption{Security goals and enforcement mechanisms.}
\label{tab:security-goals}
\begin{tabular}{p{2.2cm}p{2.8cm}p{2.2cm}}
\toprule
\textbf{Goal} & \textbf{Description} & \textbf{Mechanism} \\
\midrule
Detection Integrity & Accurate threat identification & eBPF tracing + ML \\
Rule Isolation & Malicious rules cannot harm host & Wasm sandbox \\
Data Confidentiality & Sensitive events protected & Memory isolation \\
Availability & System remains responsive & Resource limits \\
\bottomrule
\end{tabular}
\end{table}

\textbf{G1: Detection Integrity.}
The system must accurately identify malicious activities with high recall (to catch attacks) and reasonable precision (to avoid alert fatigue).
Attackers should not be able to trivially evade detection through simple obfuscation.

\textbf{G2: Rule Isolation.}
Third-party detection rules must be confined to their sandbox.
A compromised or malicious rule should not be able to:
(1) access host file system or network;
(2) read memory of other Wasm modules;
(3) interfere with the host program or kernel.

\textbf{G3: Data Confidentiality.}
Security events may contain sensitive information (file paths, network destinations, command lines).
Rules should only access the specific events relevant to their detection logic, enforced through capability-based access control.

\textbf{G4: Availability.}
The detection system must remain responsive even under attack.
Resource exhaustion attempts by malicious rules or event floods from external attackers should not render the system unresponsive.

\subsection{Trusted Computing Base}

Our TCB consists of the following components:

\begin{figure}[t]
\centering
\begin{tikzpicture}[
    box/.style={rectangle, draw, minimum width=2.5cm, minimum height=0.8cm, align=center},
    trusted/.style={box, fill=green!20},
    untrusted/.style={box, fill=red!10}
]
    % Trusted components
    \node[trusted] (kernel) at (0,0) {Linux Kernel\\+ eBPF};
    \node[trusted] (wasmtime) at (3,0) {Wasmtime\\Runtime};
    \node[trusted] (host) at (6,0) {Host Program\\(Go)};

    % Untrusted
    \node[untrusted] (rules) at (3,-1.5) {Detection Rules\\(Wasm)};

    % Labels
    \node[above=0.3cm of kernel, font=\footnotesize\bfseries] {TCB (Trusted)};
    \node[below=0.3cm of rules, font=\footnotesize\bfseries] {Untrusted};

    % Boundary
    \draw[dashed, thick] (-1.5,-0.7) -- (7.5,-0.7);
\end{tikzpicture}
\caption{Trusted computing base of \textsc{WasmSentinel}.}
\label{fig:tcb}
\end{figure}

\begin{itemize}
    \item \textbf{Linux Kernel + eBPF:} We trust the kernel's eBPF verifier to ensure safety of eBPF programs. This is a standard assumption for eBPF-based security tools.

    \item \textbf{Wasmtime Runtime:} We trust the Wasm runtime to correctly enforce sandbox isolation. Wasmtime is developed by the Bytecode Alliance and has undergone extensive security review.

    \item \textbf{Host Program:} Our Go-based host application is trusted to correctly orchestrate eBPF event collection and Wasm rule execution.
\end{itemize}

Detection rules, including the ML models, are \textit{untrusted}.
This is a key design decision that enables deployment of community-contributed rules without compromising system security.
